
    \documentclass[a4paper]{article}
    \usepackage{savetrees}
    \usepackage{hyperref}
    \usepackage{pdfpages}

    \begin{document}

    Krishna Pothugunta and John P Lalor.  {\bf \href{http://cs.umd.edu/~jbg//docs/2026_aclfindings_cct.pdf}{Carefully Considering Culture: Analyzing LLM Alignment in Single-and Multi-Cultural Settings using Cultural Consensus Theory}}.  \emph{Findings of the Association for Computational Linguistics}, 2026.



    \vspace{5cm}

    \begin{verbatim}
    @article{Pothugunta:Lalor-2026,
	Title = {Carefully Considering Culture: Analyzing LLM Alignment in Single-and Multi-Cultural Settings using Cultural Consensus Theory},
	Author = {Krishna Pothugunta and John P Lalor},
	Journal = {Findings of the Association for Computational Linguistics},
	Year = {2026},
	Url = {http://cs.umd.edu/~jbg//docs/2026_aclfindings_cct.pdf},
}

    \end{verbatim}

    {\bf Accessible Abstract:} Recent work in NLP has probed large language models for their understanding of cultural norms across countries. However, this work typically considers distributional patterns, ignoring group consensus or possible multicultural environments within a country. In this work, we leverage cultural consensus theory (CCT) from cultural anthropology to model such multidimensional nuance. Applying CCT to the World Values Survey (WVS) across 10 countries and 12 domains, we demonstrate that models frequently misrepresent cultural structures by either failing to form cohesive consensus or severely over-regularizing consensus. Through explicit representation of intra-group variance, CCT provides actionable diagnostics to evaluate when models reflect true human diversity versus algorithmic homogenization.\\~

    

    

    \vspace{3cm}
    Downloaded from \url{http://cs.umd.edu/~jbg/docs/2026_aclfindings_cct.pdf}  \\

    \vspace{2cm}
    \textit{Contact Jordan Boyd-Graber (jbg@boydgraber.org) for questions about this paper.}

    \includepdf[pages={-}]{src_docs/2026_aclfindings_cct.pdf}

    \end{document}
    